\documentclass[11pt,a4paper]{article}
\usepackage[T1]{fontenc}
\usepackage{lmodern}
\usepackage[margin=26mm]{geometry}
\usepackage{amsmath,amssymb,amsthm,booktabs,array,longtable}
\usepackage{microtype,multicol}
\usepackage[hidelinks]{hyperref}
\setlength{\parskip}{3pt}
\setlength{\emergencystretch}{2em}
\raggedbottom
\allowdisplaybreaks[1]
\clubpenalty=10000
\widowpenalty=10000
\newtheorem{theorem}{Theorem}[section]
\newtheorem{lemma}[theorem]{Lemma}
\newtheorem{corollary}[theorem]{Corollary}
\theoremstyle{remark}
\newtheorem{remark}[theorem]{Remark}
\newcommand{\Z}{\mathbb Z}
\newcommand{\Q}{\mathbb Q}
\newcommand{\N}{\mathbb N}
\newcommand{\F}{\mathbb F}
\newcommand{\ps}[1]{#1[[x]]}
\newcommand{\Rev}{\operatorname{Rev}}
\newcommand{\rev}{^{\langle-1\rangle}}
\newcommand{\Plus}{\operatorname{Plus}}
\newcommand{\Minus}{\operatorname{Minus}}
\newcommand{\coeff}[2]{[x^{#1}]#2}
\newcommand{\E}{\mathcal E}
\Urlmuskip=0mu plus 1mu
\hypersetup{pdftitle={A195194: Parity of an infinite nested reversion},pdfauthor={Dev Goyal},pdfsubject={Separate manuscript with AI disclosure and Lean proof map}}
\title{A195194\\[4pt]\large Parity of an infinite nested reversion}
\author{Dev Goyal}
\date{24 October 2026}
\begin{document}
\maketitle
\begin{center}
\fbox{\begin{minipage}{0.9\linewidth}\small
\textbf{AI disclosure.} ChatGPT (OpenAI) generated the mathematical arguments,
Lean proofs, and this manuscript in response to Dev Goyal's requests.
Dev Goyal is the sole human author. The source sequence and conjecture or
observation are credited to Paul D. Hanna.
\textbf{Independent human review is pending.}
\end{minipage}}
\end{center}
\begin{abstract}
For arbitrary weights in a commutative ring, we construct a unique infinite family of normalized nested reversions and prove finite-depth coefficient stabilization. For weights $1,2,3,\ldots$, the first tail has odd coefficients exactly at powers of two. This is a statement about formal series, independent of analytic convergence.
\end{abstract}

\section{Source, notation, and scope}

The source is Paul D. Hanna's OEIS entry A195194 \cite{oeis194},
retrieved on 24 October 2026. Its initial sequence index is 1.
We also retain the formal zero constant coefficient at index zero.

All series are formal. The set $\N$ includes zero, and $[x^n]U$ denotes a
coefficient. A normalized series is $x+O(x^2)$; $O(x^r)$ means divisibility
by $x^r$, over any coefficient ring used below. We write $U\circ V=U(V(x))$
only when $V(0)=0$. Compositional inversion, denoted $\Rev(U)$ or $U\rev$,
is distinct from multiplicative inversion of a unit. Congruences are
coefficientwise; ordinary products and compositional iterates are distinct.

\section{Formal algebra used in the proof}

\begin{lemma}[Triangular construction]\label{lem:construction}
Let $K$ be a commutative ring. Suppose $\Phi$ preserves normalized series
and improves agreement by one degree:
\[
 G-H=O(x^r)\quad\Longrightarrow\quad
 \Phi(G)-\Phi(H)=O(x^{r+1})\qquad(r\ge2).
\]
Then $\Phi$ has a unique normalized fixed point. The same conclusion holds
for a family of series if the degree-$n$ equations, at every family index,
determine the degree-$n$ coefficients from strictly lower degrees.
\end{lemma}

\begin{proof}
The degree-$n$ coefficient of $\Phi(G)$ depends only on coefficients of
$G$ below degree $n$. Fix the constant and linear coefficients, and then
define every later coefficient successively. This gives a fixed point;
the same induction proves uniqueness. For a family, perform the induction
on degree simultaneously at every family index. The same induction applies with only the constant coefficient prescribed,
provided agreement is improved starting at degree one. It also applies to
an operator on all series when agreement is improved starting at degree zero.
For the operators used in this paper, the recurrence requires
only ring operations, so it also constructs integral solutions when
$K=\Z$. Equivalently, iterating $\Phi$ from any normalized initial series
fixes one further coefficient at each step.
\end{proof}

\begin{lemma}[Inversion and transport]\label{lem:inverse}
A normalized series over any commutative ring has a unique normalized
two-sided compositional inverse. A constant-one series has a unique
multiplicative inverse with constant coefficient one. These inverses and
valid composition commute with coefficient ring homomorphisms.
Normalized compositional inversion preserves agreement modulo $x^r$.
\end{lemma}

\begin{proof}
For normalized $G$, the degree-$n$ equation in $G(H)=x$ determines $h_n$
with multiplier one from earlier coefficients. Construct a right inverse
$J$ of $H$. Associativity gives
$G=G\circ(H\circ J)=(G\circ H)\circ J=J$, so $H\circ G=x$ as well.
The same associativity argument proves uniqueness. Multiplicative
inversion follows similarly from the coefficient recurrence for $UV=1$.

Each coefficient of a valid composition is a finite polynomial in the
coefficients of its arguments. Ring homomorphisms preserve these
polynomials and the defining equations, so uniqueness gives transport of
inversion. If $G-H=O(x^r)$, composition with any normalized series
preserves that agreement. Composing with the respective inverses then
shows $\Rev(G)-\Rev(H)=O(x^r)$.
\end{proof}

\begin{lemma}[Binary series and doubling]\label{lem:binary}
There is a unique normalized integral series $C=x+C^2$. Its reduction
modulo two is
\[
 P(x)=\sum_{k\ge0}x^{2^k}.
\]
For series $U,V$ over $\Z/4\Z$, the equality $2U=2V$ is equivalent to
equality of their reductions modulo two.
\end{lemma}

\begin{proof}
The operator $G\mapsto x+G^2$ improves agreement because
$G^2-H^2=(G-H)(G+H)$ and normalized $G+H$ has zero constant coefficient.
In characteristic two, $U(x)^2=U(x^2)$ for every formal series. Thus the
normalized equation $U=x+U^2$ forces the coefficient at one to be one,
every higher odd-index coefficient to be zero, and the coefficient at
$2m$ to equal the coefficient at $m$. This gives exactly $P$.

For scalar residues modulo four, $2u=2v$ holds precisely when $u$ and $v$
have the same parity. Apply this statement coefficientwise. In
particular, this argument does not cancel a factor of two in a ring in
which two is a zero divisor.
\end{proof}

\section{A195194: infinite nested reversion}

The title of A195194 is an infinite sequence of nested reversions
\cite{oeis194}. It is useful to construct the more general family with
arbitrary weights before specializing to $1,2,3,\ldots$.

\begin{theorem}[Canonical nested tails]\label{thm:tails}
For every sequence $(w_j)_{j\ge0}$ in a commutative ring $K$, there is a
unique family of normalized series satisfying
\[
 F_j=\Rev(x-w_jxF_{j+1})\qquad(j\ge0). \tag{1}
\]
If the finite nests are defined by $F_j^{[0]}=x$ and
\[
 F_j^{[d+1]}=\Rev(x-w_jxF_{j+1}^{[d]}),
\]
then $F_j^{[d]}-F_j=O(x^{d+2})$ for every $d,j\ge0$.
\end{theorem}

\begin{proof}
The right inverse identity in (1) is equivalent to
\[
 F_j=x+w_jF_j\,F_{j+1}(F_j). \tag{2}
\]
Both factors in the nontrivial term start in degree one. Consequently its
degree-$n$ coefficient depends only on coefficients below degree $n$ of
$F_j$ and $F_{j+1}$. Lemma~\ref{lem:construction} constructs the unique
family over $K$. No boundary condition at infinity is assumed.
Since $x-w_jxF_{j+1}$ is normalized, Lemma~\ref{lem:inverse} turns the
right inverse equation into both inverse identities, proving (1).

At depth zero both the approximation and the true tail start with $x$,
so they agree modulo $x^2$. Multiplication by $x$ improves the agreement
of their next tails by one degree, and compositional inversion preserves
that improved agreement. Induction on $d$ gives the stated bound. Thus
the finite nested expressions agree with the canonical family through
degree $d+1$; their formal infinite limit is unambiguous.
\end{proof}

\begin{theorem}[A195194 parity]\label{thm:nestedparity}
Let $w_j=j+1$, and let $A=F_0$ in Theorem~\ref{thm:tails}, over $\Z$.
For every $n\in\N$,
\[
 [x^n]A\text{ is odd}\quad\Longleftrightarrow\quad
 n=2^k\text{ for some }k\in\N.
\]
\end{theorem}

\begin{proof}
Map the family to $\F_2[[x]]$ and write a bar for reduction. The second
weight is two, so (2) gives $\overline F_1=x$. The first weight is one,
and therefore
\[
 \overline A=x+\overline A\,\overline F_1(\overline A)
             =x+\overline A^2.
\]
Lemma~\ref{lem:binary} identifies $\overline A$ with $P$. The conclusion
also holds at $n=0$, where the coefficient is zero and the index is not a
power of two.
\end{proof}

The formal implementation constructs the simultaneous family using the
coefficient ring $\N\to K$, with its shift endomorphism. Evaluation at an
index recovers (2). The theorem \texttt{source\_existence\_uniqueness}
states the general source result using both inverse identities;
\texttt{approximation\_agrees} supplies the finite-depth bound, and
\texttt{a195194\_parity} gives the infinite parity statement. All are in
\texttt{Round4.NestedReversion}.


\section{Separate Lean proof and checked scope}

The entry file is \nolinkurl{lean/Entries/A195194.lean}. Its explicit
theorem statements and proof terms use the namespace
\nolinkurl{FinalCollection.A195194}. The complete imported implementations
and their shared mathematical dependencies are included in the collection.
The following local theorem names give the proof map:

\begin{multicols}{2}\raggedright\begin{itemize}\setlength{\itemsep}{2pt}\setlength{\parsep}{0pt}
\item \nolinkurl{normalized_source}.
\item \nolinkurl{all_tails_normalized}.
\item \nolinkurl{all_tails_source}.
\item \nolinkurl{nested_definition_unique}.
\item \nolinkurl{first_tail_reversion}.
\item \nolinkurl{finite_nest_coefficients}.
\item \nolinkurl{binary_equation}.
\item \nolinkurl{exact_parity}.
\item \nolinkurl{arbitrary_weights_unique}.
\item \nolinkurl{arbitrary_finite_depth}.
\end{itemize}\end{multicols}

The main implementation module is \nolinkurl{Round4.NestedReversion}.
The catalogue records the full original endpoint names and the new entry
theorems. The project pins Lean 4.19.0 and Mathlib revision
\nolinkurl{c44e0c8ee63ca166450922a373c7409c5d26b00b}. The collection
report records source recompilation and transitive axiom checks. Accepted
axioms are propositional extensionality, classical choice, and quotient
soundness. All index conclusions above are unbounded; exact finite
coefficient checks in the package provide separate corroboration.

\paragraph{Limits and status.}
Convergence means coefficient stabilization of formal series, with no analytic assertion. The source starts at n=1.
The manuscript consolidates completed results; it does not assert global
novelty, exclusive priority, peer review, OEIS acceptance, or a new submission.
Standard algebraic ingredients are not claimed as new discoveries.

\begin{thebibliography}{9}
\raggedright
\bibitem{oeis194}
Paul D. Hanna, \emph{OEIS A195194}, 11 September 2011.
\url{https://oeis.org/A195194}. Retrieved 24 October 2026.
\end{thebibliography}
\end{document}
